在前面的章节中我们已经看到，C++ 与动态内存分配之间有着一种爱恨交织的关系。

一方面，从堆（heap）进行的动态内存分配是一种「代码异味（code smell）」；追逐指针会损害程序的性能，堆可能在毫无预兆的情况下被耗尽（从而引发 \texttt{std::\allowbreak{}bad\_\allowbreak{}alloc} 类型的异常），而手动内存管理又极其微妙且困难，以至于 C++11 引入了几种不同的「智能指针（smart pointer）」类型来应对这种复杂性（见第 6 章「智能指针」）。2011 年之后陆续推出的 C++ 版本，还新增了大量不进行分配（non-allocating）的代数数据类型，例如 \texttt{tuple}、\texttt{optional} 和 \texttt{variant}（见第 5 章「词汇类型」），它们无需触碰堆即可表达所有权或包含关系。

另一方面，这些新的智能指针类型确实有效地管理了内存管理的复杂性；在现代 C++ 中，你可以安全地分配和释放内存，而无需使用原生的 \texttt{new} 或 \texttt{delete}，也不必担心内存泄漏。而且，堆分配被许多新 C++ 特性（\texttt{any}、\texttt{function}、\texttt{promise}）在「背后（under the hood）」使用，正如它继续被许多旧特性（\texttt{stable\_\allowbreak{}partition}、\texttt{vector}）使用一样。

于是这里就出现了一个矛盾：如果同时有人告诫我们，优秀的 C++ 代码应当避免堆分配，那么我们又该如何使用这些依赖堆分配的优秀新特性（以及旧特性）呢？

在大多数情况下，你应该倾向于使用 C++ 提供的特性。如果你想要一个可调整大小的元素向量，就应当使用默认的 \texttt{std::\allowbreak{}vector}，除非你已经实测在你的场景中它确实存在性能问题。但也确实存在一类程序员——他们在诸如飞行软件这样高度受限的环境中工作——由于一个非常简单的原因而必须避免触碰堆：在他们的平台上，「堆」根本不存在！在这些嵌入式环境中，整个程序的内存占用必须在编译期就安排妥当。这类程序中有一些索性避免使用任何类似堆分配的算法——如果你从不动态分配任何资源，就永远不可能遭遇意料之外的资源耗尽！另一些程序确实使用了类似堆分配的算法，但要求「堆」在程序中显式地表示出来（例如，用一个很大的 \texttt{char} 数组，再加上一些用于「预约」和「归还」该数组连续区块的函数）。

如果最后这一类程序无法使用 C++ 所提供的特性（例如 \texttt{std::\allowbreak{}vector} 和 \texttt{std::\allowbreak{}any}），那将是非常遗憾的。因此，自 1998 年的第一版标准起，标准库就提供了一种被称为「分配器感知（allocator-awareness）」的特性。

当一个类型或算法具备分配器感知能力时，它就提供了一种途径，让程序员能够精确地指定该类型或算法应当如何预约和归还动态内存。这个「如何做」被具体化为一个被称为分配器（allocator）的对象。

在本章中，我们将学习：

\begin{itemize}
\item 「分配器」与「内存资源（memory resource）」的定义
\item 如何创建一个从静态缓冲区分配的内存资源
\item 如何让自己的容器具备「分配器感知」能力
\item 命名空间 \texttt{std::\allowbreak{}pmr} 中的标准内存资源类型，以及它们令人意外的陷阱
\item C++11 分配器模型中许多奇怪的特性，纯粹是为了支持 \texttt{scoped\_\allowbreak{}allocator\_\allowbreak{}adaptor} 而设计的
\item 是什么让一个类型成为「花式指针（fancy pointer）」类型，以及这类类型可能在哪些地方派上用场
\end{itemize}

\section{分配器是内存资源的句柄}

在阅读本章时，你必须牢记两个基本概念之间的区别，我将它们分别称为内存资源（memory resource）与分配器（allocator）。内存资源（其名称受到了标准自身术语的启发——你可能觉得把它叫做「堆」更自然）是一个长寿命的对象，能够按需分发出内存块（通常是从内存资源自身所拥有的一大块内存中切分出来）。内存资源具有典型的面向对象（OO）语义（见第 1 章「古典多态与泛型编程」）：你创建一次内存资源，之后就再也不移动或拷贝它，而内存资源的相等性通常由对象身份来定义。另一方面，分配器（allocator）是一个指向内存资源的短寿命句柄。分配器具有指针语义：你可以拷贝它们、在四处移动它们，基本上随心所欲地摆弄它们，而分配器的相等性通常由它们是否指向同一个内存资源来定义。除了说一个分配器「指向」某个特定的内存资源，我们也可以说该分配器「由」那个内存资源「支撑（backed by）」；这两种说法是等价的。

当我在本章中谈到「内存资源」和「分配器」时，我指的是前面所说的那些概念。标准库中也确实有几个名为 \texttt{memory\_\allowbreak{}resource} 和 \texttt{allocator} 的类型；每当我谈论这些类型时，我会谨慎地使用打字机字体。这应该不至于造成太大困惑。这种情况与第 2 章「迭代器与范围」类似，在那章里我们既谈到了「迭代器」，也谈到了 \texttt{std::\allowbreak{}iterator}。当然，那一章要简单些，因为我提到 \texttt{std::\allowbreak{}iterator} 只是为了告诉你永远不要使用它；在写得好看的 C++ 代码中，它毫无用武之地。而在本章中我们将了解到，\texttt{std::\allowbreak{}pmr::\allowbreak{}memory\_\allowbreak{}resource} 在某些 C++ 程序里确实有它的一席之地！

尽管我把分配器描述为一个「指向」内存资源的句柄，但你应该注意到，有时候我们所讨论的内存资源是一个全局单例——这类单例的一个例子就是全局堆，它的访问函数正是全局的 \texttt{operator} \texttt{new} 和 \texttt{operator delete}。正如一个「捕获」了全局变量的 lambda 实际上并没有捕获任何东西一样，由全局堆支撑的分配器实际上也不需要任何状态。事实上，\texttt{std::\allowbreak{}allocator<\allowbreak{}T>\allowbreak{}} 正是这样一种无状态（stateless）的分配器类型——不过这里我们有些操之过急了！

\subsection{回顾——接口与概念}

回顾一下第 1 章「古典多态与泛型编程」：C++ 提供了两种基本互不相容的方式来处理多态性。静态的、编译期的多态被称为泛型编程（generic programming）；它依赖把一个多态接口表达为一个带许多可能模型的概念（concept），而与接口交互的代码则以模板的形式表达。动态的、运行期的多态被称为古典多态（classical polymorphism）；它依赖把一个多态接口表达为一个带有许多可能派生类的基类，而与接口交互的代码则以调用虚方法的形式表达。

在本章中，我们将与泛型编程进行第一次（也是最后一次）真正意义上的亲密接触。要理解 C++ 的分配器，除非你能同时在脑海中把握两个观念：一方面是一个定义了接口的概念（concept）\texttt{Allocator}，另一方面是某个特定的模型，例如实现了符合 \texttt{Allocator} 概念之行为的 \texttt{std::\allowbreak{}allocator}。

更复杂的是，\texttt{Allocator} 概念实际上是一个模板化的概念族！更准确的说法是去谈论概念族 \texttt{Allocator<\allowbreak{}T>\allowbreak{}}；例如，\texttt{Allocator<\allowbreak{}int>\allowbreak{}} 将是这样一个概念：「分配 \texttt{int} 对象的分配器」，而 \texttt{Allocator<\allowbreak{}char>\allowbreak{}} 则是「分配 \texttt{char} 对象的分配器」，依此类推。再例如，具体的类 \texttt{std::\allowbreak{}allocator<\allowbreak{}int>\allowbreak{}} 是概念 \texttt{Allocator<\allowbreak{}int>\allowbreak{}} 的一个模型，但它并不是 \texttt{Allocator<\allowbreak{}char>\allowbreak{}} 的模型。

每个 \texttt{T} 的分配器（每个 \texttt{Allocator<\allowbreak{}T>\allowbreak{}}）都必须提供一个名为 \texttt{allocate} 的成员函数，使得 \texttt{a.\allowbreak{}allocate(\allowbreak{}n)\allowbreak{}} 返回一个指针，指向足以容纳 \texttt{n} 个 \texttt{T} 类型对象的数组的内存。（该指针将来自支撑该分配器实例的内存资源。）并没有规定 \texttt{allocate} 成员函数应该是静态的还是非静态的，也没有规定它是否应当只接受一个参数（\texttt{n}），或者也可能接受一些带有默认值的额外参数。因此，下面这两个类类型在上述意义上都称得上是 \texttt{Allocator<\allowbreak{}int>\allowbreak{}} 的可接受模型：

\begin{codebox}[cpp]
    struct int_allocator_2014 {
      int *allocate(size_t n, const void *hint = nullptr);
    };
    struct int_allocator_2017 {
      int *allocate(size_t n);
    };
\end{codebox}

被命名为 \texttt{int\_\allowbreak{}allocator\_\allowbreak{}2017} 的类显然是建模 \texttt{Allocator<\allowbreak{}int>\allowbreak{}} 的一种更简单的方式，但 \texttt{int\_\allowbreak{}allocator\_\allowbreak{}2014} 同样是一个正确的模型，因为在两种情况下，表达式 \texttt{a.\allowbreak{}allocate(\allowbreak{}n)\allowbreak{}} 都会被编译器接受；而当我们谈论泛型编程时，这就是我们所要求的全部。

相反，当我们使用古典多态时，我们会为基类的每个方法规定一个固定的签名，而派生类不允许偏离该签名：

\begin{codebox}[cpp]
    struct classical_base {
      virtual int *allocate(size_t n) = 0;
    };
    struct classical_derived : public classical_base {
      int *allocate(size_t n) override;
    };
\end{codebox}

派生类 \texttt{classical\_\allowbreak{}derived} 不允许在 \texttt{allocate} 方法的签名上添加任何额外参数；不允许改变返回类型；也不允许把该方法变成 \texttt{static}。相比于泛型编程，古典多态的接口被「锁定（locked down）」得更严。

由于一个被「锁定」的古典接口，天生就比一个开放的概念性接口更容易描述，所以我们将从 C++17 崭新推出的、具备古典多态性的 \texttt{memory\_\allowbreak{}resource} 开始，来游览分配器库。

\section{用 \texttt{memory\_\allowbreak{}resource} 定义堆}

请回想一下：在资源受限的平台上，我们可能不被允许使用「堆」（例如通过 \texttt{new} 和 \texttt{delete}），因为平台的运行时可能不支持动态内存分配。但我们完全可以自己造一个小堆——不是「那个堆」，而只是「一个堆」——并通过编写几个函数 \texttt{allocate} 和 \texttt{deallocate} 来模拟动态内存分配的效果，这两个函数从一大块静态分配的 \texttt{char} 数组中预约出区块，大致如下：

\begin{codebox}[cpp]
    static char big_buffer[10000];
    static size_t index = 0;
    void *allocate(size_t bytes) {
      if (bytes > sizeof big_buffer - index) {
        throw std::bad_alloc();
      }
      index += bytes;
      return &big_buffer[index - bytes];
    }
    void deallocate(void *p, size_t bytes) {
      // 直接丢弃
    }
\end{codebox}

为了让代码尽量简单，我把 \texttt{deallocate} 做成了一个空操作（no-op）。这个小堆允许调用者分配至多 10,000 字节的内存，之后便开始抛出 \texttt{bad\_\allowbreak{}alloc}。

只要在代码上再多花一点功夫，我们就能让调用者无限次地进行分配与释放，前提是已分配内存的未归还总量不超过 10,000 字节，并且调用者始终遵循「最后分配者最先释放（last-allocated-first-deallocated）」的协议：

\begin{codebox}[cpp]
    void deallocate(void *p, size_t bytes) {
      if ((char*)p + bytes == &big_buffer[index]) {
        // 啊哈！我们可以回退索引了！
        // 直接丢弃
      } else {
        // 直接丢弃
      }
    }
\end{codebox}

这里的关键点是，我们的堆拥有一些状态（在本例中就是 \texttt{big\_\allowbreak{}buffer} 和 \texttt{index}），以及几个操作这些状态的函数。我们已经看到了 \texttt{deallocate} 两种可能的实现——也还存在其他带有额外共享状态的实现，它们不会这么「leaky（易漏）」——然而接口、也就是 \texttt{allocate} 和 \texttt{deallocate} 自身的签名，始终保持不变。这表明我们可以把自己的状态和访问函数封装进一个 C++ 对象；而实现方式的丰富多样，加上函数签名的恒定不变，又表明我们可以使用某种古典多态。

C++17 的分配器模型正是这样做的。标准库提供了一个具备古典多态性的基类 \texttt{std::\allowbreak{}pmr::\allowbreak{}memory\_\allowbreak{}resource} 的定义，然后我们把自己的小堆实现为一个派生类。（在实践中，我们或许会使用标准库提供的某个派生类，但让我们先把这个小例子讲完，再讨论它们。）基类 \texttt{std::\allowbreak{}pmr::\allowbreak{}memory\_\allowbreak{}resource} 定义于标准头文件 \texttt{<memory\_\allowbreak{}resource>\allowbreak{}} 中：

\begin{codebox}[cpp]
    class memory_resource {
      virtual void *do_allocate(size_t bytes, size_t align) = 0;
      virtual void do_deallocate(void *p, size_t bytes, size_t align) = 0;
      virtual bool do_is_equal(const memory_resource& rhs) const = 0;
    public:
      void *allocate(size_t bytes, size_t align) {
        return do_allocate(bytes, align);
      }
      void deallocate(void *p, size_t bytes, size_t align) {
        return do_deallocate(p, bytes, align);
      }
      bool is_equal(const memory_resource& rhs) const {
        return do_is_equal(rhs);
      }
    };
\end{codebox}

请注意这个类在 \texttt{public} 接口与 \texttt{virtual} 实现之间那层奇特的间接层（indirection）。通常当我们进行古典多态时，只有一套方法，它们既 \texttt{public} 又 \texttt{virtual}；但在此处，我们有一个 \texttt{public} 的非虚接口，它会向下调用私有的虚方法。这种把接口与实现分离开来的做法有几点不明显的益处——例如，它能阻止任何子类用「非虚地直接调用虚方法」的语法来调用 \texttt{this->Some\allowbreak{}Base\allowbreak{}Class::\allowbreak{}allocate(\allowbreak{})\allowbreak{}}——但说实话，它对我们最大的好处是：当我们定义派生类时，根本不需要使用 \texttt{public} 关键字。因为我们只指定实现、而非接口，所以所写的所有代码都可以是 \texttt{private} 的。下面是我们的那个微不足道的、会漏内存的小堆：

\begin{codebox}[cpp]
    class example_resource : public std::pmr::memory_resource {
      alignas(std::max_align_t) char big_buffer[10000];
      size_t index = 0;
      void *do_allocate(size_t bytes, size_t align) override {
        if (align > alignof(std::max_align_t) ||
            (-index % align) > sizeof big_buffer - index ||
            bytes > sizeof big_buffer - index - (-index % align))
        {
            throw std::bad_alloc();
        }
        index += (-index % align) + bytes;
        return &big_buffer[index - bytes];
      }
      void do_deallocate(void *, size_t, size_t) override {
        // 直接丢弃
      }
      bool do_is_equal(const memory_resource& rhs) const override {
        return this == &rhs;
      }
    };
\end{codebox}

请注意，标准库的 \texttt{std::\allowbreak{}pmr::\allowbreak{}memory\_\allowbreak{}resource::\allowbreak{}allocate} 不仅接受一个以字节为单位的大小，还接受一个对齐（alignment）要求。我们需要确保从 \texttt{do\_\allowbreak{}allocate} 返回的任何指针都得到了恰当的对齐；例如，如果调用者打算在我们给它的内存里存放 \texttt{int}，他可能会要求四字节对齐。

关于我们的派生类 \texttt{example\_\allowbreak{}resource}，最后要注意的一点是：它代表了由我们的「堆」所控制的实际资源；也就是说，它实实在在地包含、拥有并管理着它从中分配内存的 \texttt{big\_\allowbreak{}buffer}。对于任意给定的 \texttt{big\_\allowbreak{}buffer}，我们的程序中都会有且仅有一个 \texttt{example\_\allowbreak{}resource} 对象在操作这块缓冲区。正如我们前面所说：\texttt{example\_\allowbreak{}resource} 类型的对象是「内存资源」，因此它们并不打算被拷贝或在各处移动；它们是典型的面向对象（OO）式的，而不是值语义（value-semantic）的。

标准库提供了几种（species）内存资源，全都派生自 \texttt{std::\allowbreak{}pmr::\allowbreak{}memory\_\allowbreak{}resource}。让我们来看其中的几个。

\section{使用标准内存资源}

标准库中的内存资源有两种风味（flavors）。其中一些是真正的类类型，你可以创建它们的实例；而另一些则是「匿名」类类型，只能通过单例函数来访问。一般来说，你可以通过思考「该类型的两个对象是否有可能真的是『不同的』」，或者「该类型本质上是否就是个单例」，来预测它属于哪一种。

在 \texttt{<memory\_\allowbreak{}resource>\allowbreak{}} 头文件中，最简单的内存资源就是通过 \texttt{std::\allowbreak{}pmr::\allowbreak{}null\_\allowbreak{}memory\_\allowbreak{}resource(\allowbreak{})\allowbreak{}} 访问的那个「匿名」单例。这个函数的大致定义如下：

\begin{codebox}[cpp]
    class UNKNOWN : public std::pmr::memory_resource {
      void *do_allocate(size_t, size_t) override {
        throw std::bad_alloc();
      }
      void do_deallocate(void *, size_t, size_t) override {}
      bool do_is_equal(const memory_resource& rhs) const override {
        return this == &rhs;
      }
    };
    std::pmr::memory_resource *null_memory_resource() noexcept {
      static UNKNOWN singleton;
      return &singleton;
    }
\end{codebox}

请注意，该函数返回一个指向单例实例的指针。一般来说，\texttt{std::\allowbreak{}pmr::\allowbreak{}memory\_\allowbreak{}resource} 对象是通过指针来操作的，因为 \texttt{memory\_\allowbreak{}resource} 对象本身无法四处移动。

\texttt{null\_\allowbreak{}memory\_\allowbreak{}resource} 看起来相当没用；它所做的一切就是在你试图从它分配内存时抛出一个异常。然而，当你开始使用那些更复杂的内存资源（我们稍后就会看到）时，它却能派上用场。

下一个较为复杂的内存资源，是通过 \texttt{std::\allowbreak{}pmr::\allowbreak{}new\_\allowbreak{}delete\_\allowbreak{}resource(\allowbreak{})\allowbreak{}} 访问的单例；它使用 \texttt{::\allowbreak{}operator new} 和 \texttt{::\allowbreak{}operator delete} 来分配和释放内存。

现在我们转到命名类类型（named class types）上。这类资源之所以这样命名，是因为在一个程序中拥有多个同类型的资源是合理的。例如，有 \texttt{class std::\allowbreak{}pmr::\allowbreak{}monotonic\_\allowbreak{}buffer\_\allowbreak{}resource}。这个内存资源本质上与我们前面的 \texttt{example\_\allowbreak{}resource} 相同，只有两点区别：它不像成员数据那样持有自己的大缓冲区（\texttt{std::\allowbreak{}array} 风格），而只是持有一个指针

指向从别处分配而来的一个大缓冲区（\texttt{std::\allowbreak{}vector} 风格）。当它的第一个大缓冲区耗尽时，它不会立刻开始抛出 \texttt{bad\_\allowbreak{}alloc}，而是会尝试分配第二个大缓冲区，并从这个缓冲区中切出块来分配，直到它也被用光；此时它会分配第三个大缓冲区……如此往复，直到最终它连更多的大缓冲区都分配不出来。就像我们的 \texttt{example\_\allowbreak{}resource} 一样，在资源对象自身被销毁之前，没有任何已释放的内存会被真正回收。这里有一个有用的「泄压阀」：如果你调用方法 \texttt{a.\allowbreak{}release(\allowbreak{})\allowbreak{}}，\texttt{monotonic\_\allowbreak{}buffer\_\allowbreak{}resource} 就会释放它当前持有的所有缓冲区，有点像对向量调用 \texttt{clear(\allowbreak{})\allowbreak{}}。

当你构造一个类型为 \texttt{std::\allowbreak{}pmr::\allowbreak{}monotonic\_\allowbreak{}buffer\_\allowbreak{}resource} 的资源时，你需要告诉它两件事：它的第一个大缓冲区位于何处？以及，当那个缓冲区耗尽时，它该向谁申请另一个缓冲区？第一个问题通过提供一对参数 \texttt{void*,\allowbreak{} size\_\allowbreak{}t} 来描述第一个大缓冲区（可传 \texttt{nullptr}）来解答；第二个问题则通过提供一个指向该资源的「上游（upstream）」资源的 \texttt{std::\allowbreak{}pmr::\allowbreak{}memory\_\allowbreak{}resource*} 来解答。为「上游」资源传入一个合理的选择是 \texttt{std::\allowbreak{}pmr::\allowbreak{}new\_\allowbreak{}delete\_\allowbreak{}resource(\allowbreak{})\allowbreak{}}，从而用 \texttt{::\allowbreak{}operator} \texttt{new} 来分配新缓冲区。另一个合理的选择是传入 \texttt{std::\allowbreak{}pmr::\allowbreak{}null\_\allowbreak{}memory\_\allowbreak{}resource(\allowbreak{})\allowbreak{}}，从而为这一特定资源的内存用量设定一个硬上限。下面是后者的一个例子：

\begin{codebox}[cpp]
     alignas(16) char big_buffer[10000];
     std::pmr::monotonic_buffer_resource a(
       big_buffer, sizeof big_buffer,
       std::pmr::null_memory_resource()
     );
     void *p1 = a.allocate(100);
     assert(p1 == big_buffer + 0);
     void *p2 = a.allocate(100, 16); // 对齐
     assert(p1 == big_buffer + 112);
     // 现在释放到目前为止分配的所有内存，并重新开始。
     a.release();
     void *p3 = a.allocate(100);
     assert(p3 == big_buffer + 0);
     // 当缓冲区耗尽时，a 会向上游
     // 寻找更多缓冲区……但什么也找不到。
     try {
       a.allocate(9901);
     } catch (const std::bad_alloc&) {
       puts("The null_memory_resource did its job!");
     }
\end{codebox}

如果你忘了某个 \texttt{monotonic\_\allowbreak{}buffer\_\allowbreak{}resource} 正在使用哪个上游资源，随时可以通过调用 \texttt{a.\allowbreak{}upstream\_\allowbreak{}resource(\allowbreak{})\allowbreak{}} 来查明；该方法会返回指向构造时提供的那个上游资源的指针。

\subsection{从池资源分配}

C++17 标准库提供的最后一种内存资源（memory resource）就是我们所说的「池资源（pool resource）」。池资源管理的并非仅仅一个大缓冲区（例如 \texttt{example\_\allowbreak{}resource}），甚至连一条单调递增的缓冲区链（例如 \texttt{monotonic\_\allowbreak{}buffer\_\allowbreak{}resource}）都不止。相反，它管理着大量各种尺寸的「块（block）」。所有同一尺寸的块都一起存放在一个「池（pool）」中，于是我们可以谈论「尺寸为 4 的块之池」「尺寸为 16 的块之池」，依此类推。当一个大小为 k 的分配请求到来时，池资源会去尺寸为 k 的块之池里找，取出一块并返回。如果尺寸为 k 的池空了，那么池资源就会尝试从它的上游资源那里再分配一些块。此外，如果到来的分配请求大到我们根本就没有对应尺寸的池，那么池资源就被允许把该请求直接转交给它的上游资源。池资源有两种：同步（synchronized）与非同步（unsynchronized），也就是说，线程安全与线程不安全。如果你打算从两个不同的线程并发地访问一个池，那么你应该使用 \texttt{std::\allowbreak{}pmr::\allowbreak{}synchronized\_\allowbreak{}pool\_\allowbreak{}resource}；而如果你确定绝不会那样做、并且追求极致速度，那么你应该使用 \texttt{std::\allowbreak{}pmr::\allowbreak{}unsynchronized\_\allowbreak{}pool\_\allowbreak{}resource}。（顺便一提，\texttt{std::\allowbreak{}pmr::\allowbreak{}monotonic\_\allowbreak{}buffer\_\allowbreak{}resource} 始终是线程不安全的；而 \texttt{new\_\allowbreak{}delete\_\allowbreak{}resource(\allowbreak{})\allowbreak{}} 实际上是线程安全的，因为它所做的无非是调用 \texttt{new} 与 \texttt{delete}。）

当你构造一个类型为 \texttt{std::\allowbreak{}pmr::\allowbreak{}synchronized\_\allowbreak{}pool\_\allowbreak{}resource} 的资源时，你需要告诉它三件事：它的各个池里应当保留哪些块尺寸；当它去上游资源那里获取更多块时，应当把多少块打包成一个「块组（chunk）」；以及谁是它的上游资源。遗憾的是，标准接口在这一块相当不尽如人意——以至于坦白讲，我建议：如果这些参数对你真的很重要，你应该自己动手实现一个派生自 \texttt{memory\_\allowbreak{}resource} 的类，而完全不要去碰标准库提供的版本。表达这些选项的语法也相当古怪：

\begin{codebox}[cpp]
     std::pmr::pool_options options;
     options.max_blocks_per_chunk = 100;
     options.largest_required_pool_block = 256;
     std::pmr::synchronized_pool_resource a(
       options,
       std::pmr::new_delete_resource()
     );
\end{codebox}

请注意，这里没有任何办法来精确指定你想要的块尺寸；这完全交由各厂商对 \texttt{synchronized\_\allowbreak{}pool\_\allowbreak{}resource} 的具体实现来决定。运气好的话，它会选择出与你的使用场景相匹配的、还不错的块尺寸；但就我个人而言，我不会依赖于这个假设。还要注意，这里也没有办法为不同的块尺寸使用不同的上游资源，更没有为「回退（fallback）」资源（即调用方请求异常尺寸分配时所用）指定不同上游资源的途径。

简而言之，在可预见的未来，我会避开标准库内建的 \texttt{pool\_\allowbreak{}resource} 派生类。但是，从 \texttt{memory\_\allowbreak{}resource} 派生出你自己的类这一基本思路是站得住脚的。如果你关心内存分配、并想管理自己的小堆，我会建议把 \texttt{memory\_\allowbreak{}resource} 引入你的代码库。

至此，我们一直只是在谈论各种分配策略，它们由不同的 \texttt{memory\_\allowbreak{}resource} 派生类所「人格化」。我们还需要看看，如何把 \texttt{memory\_\allowbreak{}resource} 挂接进标准模板库（STL）的算法与容器之中。而要做到这一点，我们就必须从 \texttt{memory\_\allowbreak{}resource} 所代表的古典多态（classical polymorphism）世界，回退到 C++03 STL 的值语义（value-semantic）世界。

\section{标准分配器的五百顶帽子}

标准分配器（allocator）模型在 2011 年想必显得惊为天人。我们马上就会看到，仅靠一个 C++ 类型，我们便能完成下面所有这些壮举：

指定一个用于分配内存的内存资源。

为每个已分配的指针标注一些元数据，这些元数据将在其整个生命周期中被一路携带，直到释放之时。

把一个容器对象与某个特定的内存资源关联起来，并确保这种关联是「粘性的（sticky）」——这个容器对象将始终使用给定的堆来进行它的分配。

把一个容器值与某个特定的内存资源关联起来，这意味着该容器可以利用值语义被高效地搬移，而不会忘记如何释放其内容。

在上述两种互斥的行为之间做出选择。

为多级容器（例如向量的向量）的各个层级指定内存分配策略。

重新定义「构造」容器内容意味着什么，例如，\texttt{vector<\allowbreak{}int>\allowbreak{}::\allowbreak{}resize} 可以被定义为对新元素做默认初始化，而不是零初始化。

要让任何一个类类型去戴这么多的帽子，简直是疯狂之举——这严重违反了单一职责原则（Single Responsibility Principle）。然而，这正是标准分配器模型所做的事情；所以，让我们来逐一解释这些特性。

请记住，「标准分配器」不过是任意一个对某个类型 \texttt{T} 满足 \texttt{Allocator<\allowbreak{}T>\allowbreak{}} 概念的类类型。标准库提供了三种标准分配器类型：\texttt{std::\allowbreak{}allocator<\allowbreak{}T>\allowbreak{}}、\texttt{std::\allowbreak{}pmr::\allowbreak{}polymorphic\_\allowbreak{}allocator<\allowbreak{}T>\allowbreak{}}，以及 \texttt{std::\allowbreak{}scoped\_\allowbreak{}allocator\_\allowbreak{}adaptor<\allowbreak{}A...>}。

让我们先来看看 \texttt{std::\allowbreak{}allocator<\allowbreak{}T>\allowbreak{}}：

\begin{codebox}[cpp]
     template<class T>
     struct allocator {
       using value_type = T;
       T *allocate(size_t n) {
         return static_cast<T *>(::operator new(n * sizeof (T)));
       }
       void deallocate(T *p, size_t) {
         ::operator delete(static_cast<void *>(p));
       }
       // NOTE 1
       template<class U>
       explicit allocator(const allocator<U>&) noexcept {}
       // NOTE 2
       allocator() = default;
       allocator(const allocator&) = default;
     };
\end{codebox}

\texttt{std::\allowbreak{}allocator<\allowbreak{}T>\allowbreak{}} 拥有 \texttt{Allocator<\allowbreak{}T>\allowbreak{}} 概念所要求的成员函数 \texttt{allocate} 与 \texttt{deallocate}。请记住，我们现在身处基于概念的泛型编程世界！古典多态的 \texttt{memory\_\allowbreak{}resource} 同样拥有名为 \texttt{allocate} 与 \texttt{deallocate} 的成员函数，但它们始终返回 \texttt{void*}，而非 \texttt{T*}。（此外，\texttt{memory\_\allowbreak{}resource::\allowbreak{}allocate(\allowbreak{})\allowbreak{}} 接受两个参数——\texttt{bytes} 与 \texttt{align}——而 \texttt{allocator<\allowbreak{}T>\allowbreak{}::\allowbreak{}allocate(\allowbreak{})\allowbreak{}} 只接受一个参数。其中的第一个原因是，\texttt{allocator<\allowbreak{}T>\allowbreak{}} 的出现早于「对齐（alignment）是个大问题」这一主流认知；请记住，\texttt{sizeof} 运算符是从 1980 年代的 C 继承来的，而 \texttt{alignof} 运算符直到 C++11 才出现。第二个原因是，在 \texttt{std::\allowbreak{}allocator<\allowbreak{}T>\allowbreak{}} 的语境下，我们知道被分配对象的类型是 \texttt{T}，因此所请求的对齐必然就是 \texttt{alignof(\allowbreak{}T)\allowbreak{}}。\texttt{std::\allowbreak{}allocator<\allowbreak{}T>\allowbreak{}} 并没有利用这一信息，因为它早于 \texttt{alignof}；但原则上它本可以利用，而这正是 \texttt{Allocator<\allowbreak{}T>\allowbreak{}} 概念只要求签名 \texttt{a.\allowbreak{}allocate(\allowbreak{}n)\allowbreak{}}、而不是 \texttt{a.\allowbreak{}allocate(\allowbreak{}n,\allowbreak{} align)\allowbreak{}} 的原因。）

标有 \texttt{NOTE 1} 的构造函数很重要；每个分配器都需要一个仿照它而写的模板化构造函数。紧随 \texttt{NOTE 2} 这一行之后的那些构造函数则无关紧要；我们在代码中显式写出它们的唯一原因是：如果不写，由于存在用户定义的构造函数（即 \texttt{NOTE 1} 那个构造函数），它们会被隐式删除。

任何标准分配器的理念在于，我们可以把它插入（plug in）作为任意标准容器的最后一个模板类型参数（第 4 章「容器动物园」），于是该容器在任何需要分配内存的时候，都会使用那个分配器，而非它惯常的机制。我们来看一个例子：

\begin{codebox}[cpp]
     template<class T>
     struct helloworld {
       using value_type = T;
       T *allocate(size_t n) {
         printf("hello world %zu\n", n);
         return static_cast<T *>(::operator new(n * sizeof (T)));
       }
       void deallocate(T *p, size_t) {
         ::operator delete(static_cast<void *>(p));
       }
     };
     void test() {
       std::vector<int, helloworld<int>> v;
       v.push_back(42); // 打印 "hello world 1"
       v.push_back(42); // 打印 "hello world 2"
       v.push_back(42); // 打印 "hello world 4"
     }
\end{codebox}

在这里，我们的类 \texttt{helloworld<\allowbreak{}int>\allowbreak{}} 建模了 \texttt{Allocator<\allowbreak{}int>\allowbreak{}}；但我们省略了那个模板化构造函数。如果我们只跟 \texttt{vector} 打交道，这没有问题，因为 \texttt{vector} 只会分配其元素类型的数组。不过，且看当我们把测试用例改成使用 \texttt{list} 时会发生什么：

\begin{codebox}[cpp]
     void test() {
       std::list<int, helloworld<int>> v;
       v.push_back(42);
     }
\end{codebox}

在 libc++ 下，这段代码会喷吐出几十条错误信息，归结起来就是一句核心抱怨：「无法从 \texttt{helloworld<\allowbreak{}int>\allowbreak{}} 转换到 \texttt{helloworld<\allowbreak{}std::\allowbreak{}\_\allowbreak{}\_\allowbreak{}1::\allowbreak{}\_\allowbreak{}\_\allowbreak{}list\_\allowbreak{}node<\allowbreak{}int,\allowbreak{} void *>>}。」请回忆第 4 章「容器动物园」中的那张图：\texttt{std::\allowbreak{}list<\allowbreak{}T>\allowbreak{}} 把它的元素存放在比 \texttt{T} 本身更大的节点里。因此，\texttt{std::\allowbreak{}list<\allowbreak{}T>\allowbreak{}} 并不会去尝试分配任何 \texttt{T} 对象；它想要分配的是类型为 \texttt{\_\allowbreak{}\_\allowbreak{}list\_\allowbreak{}node} 的对象。要为 \texttt{\_\allowbreak{}\_\allowbreak{}list\_\allowbreak{}node} 对象分配内存，它需要一个建模了 \texttt{Allocator<\_\allowbreak{}\_\allowbreak{}list\_\allowbreak{}node>\allowbreak{}} 概念、而非 \texttt{Allocator<\allowbreak{}int>\allowbreak{}} 的分配器。

在内部，\texttt{std::\allowbreak{}list<\allowbreak{}int>\allowbreak{}} 的构造函数会接收我们的 \texttt{helloworld<\allowbreak{}int>\allowbreak{}}，并尝试把它「重绑定（rebind）」为分配 \texttt{\_\allowbreak{}\_\allowbreak{}list\_\allowbreak{}node} 对象、而非 \texttt{int} 对象。这是借助一个特性类（traits class）来完成的——这是我们曾在第 2 章「迭代器与范围」中首次遇到的一种 C++ 惯用法：

\begin{codebox}[cpp]
     using AllocOfInt = helloworld<int>;
     using AllocOfChar =
       std::allocator_traits<AllocOfInt>::rebind_alloc<char>;
     // 现在 alloc_of_char 就是 helloworld<char>
\end{codebox}

标准类模板 \texttt{std::\allowbreak{}allocator\_\allowbreak{}traits<\allowbreak{}A>\allowbreak{}} 把关于分配器类型 \texttt{A} 的大量信息整合到一处，因而易于取用。举例来说，\texttt{std::\allowbreak{}allocator\_\allowbreak{}traits<\allowbreak{}A>\allowbreak{}::\allowbreak{}value\_\allowbreak{}type} 是类型 \texttt{T} 的别名，其中 \texttt{T} 正是由 \texttt{A} 分配内存的那个类型；而 \texttt{std::\allowbreak{}allocator\_\allowbreak{}traits<\allowbreak{}A>\allowbreak{}::\allowbreak{}pointer} 则是相应指针类型的别名（通常是 \texttt{T*}）。嵌套的别名模板 \texttt{std::\allowbreak{}allocator\_\allowbreak{}traits<\allowbreak{}A>\allowbreak{}::\allowbreak{}rebind\_\allowbreak{}alloc<\allowbreak{}U>\allowbreak{}} 是一种把一个分配器从某个类型 \texttt{T}「转换」到另一个类型 \texttt{U} 的手段。这个类型特性（type trait）借助元编程撬开类型 \texttt{A}，去看：第一，\texttt{A} 是否拥有一个嵌套的模板别名 \texttt{A::\allowbreak{}rebind<\allowbreak{}U>\allowbreak{}::\allowbreak{}other}（这很少见）；第二，类型 \texttt{A} 能否被表达成 \texttt{Foo<\allowbreak{}Bar,\allowbreak{}Baz...>} 的形式（其中 \texttt{Baz...} 是某个类型列表，可能为空）。如果 \texttt{A} 可以这样表达，那么 \texttt{std::\allowbreak{}allocator\_\allowbreak{}traits<\allowbreak{}A>\allowbreak{}::\allowbreak{}rebind\_\allowbreak{}alloc<\allowbreak{}U>\allowbreak{}} 就会成为 \texttt{Foo<\allowbreak{}U,\allowbreak{}Baz...>} 的同义词。从哲学上讲，这完全是任意的；但在实践中，它对你会见到的每一种分配器类型都有效。尤其是，它对 \texttt{helloworld<\allowbreak{}int>\allowbreak{}} 也有效——这正解释了为什么我们不必在 \texttt{helloworld} 类里费心提供一个嵌套别名 \texttt{rebind<\allowbreak{}U>\allowbreak{}::\allowbreak{}other}。通过提供一个合理的默认行为，\texttt{std::\allowbreak{}allocator\_\allowbreak{}traits} 模板替我们省去了一部分样板代码。这就是 \texttt{std::\allowbreak{}allocator\_\allowbreak{}traits} 存在的理由。

你或许会奇怪，为什么 \texttt{std::\allowbreak{}allocator\_\allowbreak{}traits<\allowbreak{}Foo<\allowbreak{}Bar,\allowbreak{}Baz...>>::\allowbreak{}value\_\allowbreak{}type} 不默认就是 \texttt{Bar}。坦白讲，这我也不知道。这看起来像是板上钉钉的事；但标准库并没有这么做。因此，你编写的每一个分配器类型（请注意，现在我们谈论的是建模了 \texttt{Allocator<\allowbreak{}T>\allowbreak{}} 的类，而不是派生自 \texttt{memory\_\allowbreak{}resource} 的类）都必须提供一个嵌套的 typedef \texttt{value\_\allowbreak{}type}，它是 \texttt{T} 的别名。

不过，一旦你为 \texttt{value\_\allowbreak{}type} 定义了嵌套的 typedef，你就可以依赖 \texttt{std::\allowbreak{}allocator\_\allowbreak{}traits} 去推断出它那些嵌套 typedef 的正确定义：\texttt{pointer}（即 \texttt{T*}）、\texttt{const\_\allowbreak{}pointer}（即 \texttt{const T*}）、\texttt{void\_\allowbreak{}pointer}（即 \texttt{void*}），等等。如果你一直跟得上前面关于 \texttt{rebind\_\allowbreak{}alloc} 的讨论，你或许会猜到：把一个指针类型（如 \texttt{T*}）「转换」成 \texttt{void*}，与把一个分配器类型 \texttt{Foo<\allowbreak{}T>\allowbreak{}}「转换」成 \texttt{Foo<\allowbreak{}void>\allowbreak{}} 一样困难或一样容易；而你猜得没错！这些与指针相关的类型别名的值，全都是经由第二个标准特性类 \texttt{std::\allowbreak{}pointer\_\allowbreak{}traits<\allowbreak{}P>\allowbreak{}} 计算出来的：

\begin{codebox}[cpp]
     using PtrToInt = int*;
     using PtrToChar =
       std::pointer_traits<PtrToInt>::rebind<char>;
     // 现在 PtrToChar 是 char*
     using PtrToConstVoid =
       std::pointer_traits<PtrToInt>::rebind<const void>;
     // 现在 PtrToConstVoid 是 const void*
\end{codebox}

当我们谈到 \texttt{Allocator<\allowbreak{}T>\allowbreak{}} 的下一个职责——也就是「为每个已分配的指针标注一些元数据，这些元数据将在其整个生命周期中都被携带」——时，这个特性类就会变得非常重要。

\subsection{用花式指针（fancy pointer）携带元数据}

请考虑下面这个内存资源（memory resource）的高层设计，它应当会让你联想到 \texttt{std::\allowbreak{}pmr::\allowbreak{}monotonic\_\allowbreak{}buffer\_\allowbreak{}resource}：

维护一张我们从系统处获取到的内存块（chunk）清单。对每一块，还要记录一个 \texttt{index}，表示我们从该块起始处已分配了多少字节；并设置一个计数 \texttt{freed}，表示我们从这一特定块中已释放了多少字节。

当有人调用 \texttt{allocate(\allowbreak{}n)\allowbreak{}} 时，如果可能，就把我们任一区块的 \texttt{index} 增加相应的字节数；若实在必要，则从上游资源那里获取一个新的块。

当有人调用 \texttt{deallocate(\allowbreak{}p,\allowbreak{} n)\allowbreak{}} 时，先弄清楚 \texttt{p} 来自我们的哪一个块，再将其 \texttt{freed += n}。如果 \texttt{freed == index}，那么整块就空了，于是把 \texttt{freed = index = 0}。

把上述描述转写成代码相当直截了当。唯一成问题的一点是：在 \texttt{deallocate(\allowbreak{}p,\allowbreak{} n)\allowbreak{}} 中，我们如何判断 \texttt{p} 来自哪一个块？

\begin{codebox}[cpp]
     template<class T>
     class ChunkyPtr {
       T *m_ptr = nullptr;
       Chunk *m_chunk = nullptr;
     public:
       explicit ChunkyPtr(T *p, Chunk *ch) :
       m_ptr(p), m_chunk(ch) {}
       T& operator *() const {
         return *m_ptr;
       }
       explicit operator T *() const {
         return m_ptr;
       }
       // ... and so on ...
       // ……再加上这个额外的访问函数：
       auto chunk() const {
         return m_chunk;
       }
     };
\end{codebox}

这样一来，在我们的 \texttt{deallocate(\allowbreak{}p,\allowbreak{} n)\allowbreak{}} 函数里，我们要做的全部事情就只是查看 \texttt{p.\allowbreak{}chunk(\allowbreak{})\allowbreak{}}。但为了让这奏效，我们需要修改 \texttt{allocate(\allowbreak{}n)\allowbreak{}} 与 \texttt{deallocate(\allowbreak{}p,\allowbreak{} n)\allowbreak{}} 这两个函数的签名，使得 \texttt{deallocate} 接受 \texttt{Chunky\allowbreak{}Ptr<\allowbreak{}T>\allowbreak{}} 而非 \texttt{T*}，并且 \texttt{allocate} 返回 \texttt{Chunky\allowbreak{}Ptr<\allowbreak{}T>\allowbreak{}} 而非 \texttt{T*}。

求值为 \texttt{Chunky\allowbreak{}Ptr<\allowbreak{}T>\allowbreak{}}：

\begin{codebox}[cpp]
    template<class T>
    struct ChunkyAllocator {
      using value_type = T;
      using pointer = ChunkyPtr<T>;
      ChunkyAllocator(ChunkyMemoryResource *mr) :
        m_resource(mr) {}
      template<class U>
      ChunkyAllocator(const ChunkyAllocator<U>& rhs) :
        m_resource(rhs.m_resource) {}
      pointer allocate(size_t n) {
        return m_resource->allocate(
          n * sizeof(T), alignof(T));
      }
      void deallocate(pointer p, size_t n) {
        m_resource->deallocate(
          p, n * sizeof(T), alignof(T));
      }
    private:
      ChunkyMemoryResource *m_resource;
      template<class U>
      friend struct ChunkyAllocator;
    };
\end{codebox}

特征类 \texttt{std::\allowbreak{}allocator\_\allowbreak{}traits} 和 \texttt{std::\allowbreak{}pointer\_\allowbreak{}traits} 会负责推导其余的类型定义——例如 \texttt{void\_\allowbreak{}pointer}，它通过 \texttt{pointer\_\allowbreak{}traits::\allowbreak{}rebind} 的魔法，最终成为 \texttt{Chunky\allowbreak{}Ptr<\allowbreak{}void>\allowbreak{}} 的别名。

这里我略去了 \texttt{allocate} 和 \texttt{deallocate} 函数的实现，因为它们的具体形式取决于 \texttt{Chunky\allowbreak{}Memory\allowbreak{}Resource} 的接口。我们不妨把 \texttt{Chunky\allowbreak{}Memory\allowbreak{}Resource} 实现成下面这样：

\begin{codebox}[cpp]
    class Chunk {
      char buffer[10000];
      int index = 0;
      int freed = 0;
    public:
      bool can_allocate(size_t bytes) {
        return (sizeof buffer - index) >= bytes;
      }
      auto allocate(size_t bytes) {
        index += bytes;
        void *p = &buffer[index - bytes];
        return ChunkyPtr<void>(p, this);
      }
      void deallocate(void *, size_t bytes) {
        freed += bytes;
        if (freed == index) {
            index = freed = 0;
        }
      }
    };
    class ChunkyMemoryResource {
      std::list<Chunk> m_chunks;
    public:
      ChunkyPtr<void> allocate(size_t bytes, size_t align) {
        assert(align <= alignof(std::max_align_t));
        bytes += -bytes % alignof(std::max_align_t);
        assert(bytes <= 10000);
        for (auto&& ch : m_chunks) {
          if (ch.can_allocate(bytes)) {
            return ch.allocate(bytes);
          }
        }
        return m_chunks.emplace_back().allocate(bytes);
      }
      void deallocate(ChunkyPtr<void> p, size_t bytes, size_t) {
        bytes += -bytes % alignof(std::max_align_t);
        p.chunk()->deallocate(static_cast<void*>(p), bytes);
      }
    };
\end{codebox}

现在，我们可以用 \texttt{Chunky\allowbreak{}Memory\allowbreak{}Resource} 为标准的支持分配器的（allocator-aware）容器分配内存，就像这样：

\begin{codebox}[cpp]
    ChunkyMemoryResource mr;
    std::vector<int, ChunkyAllocator<int>> v{&mr};
    v.push_back(42);
    // v 底层数组所用的全部内存
    // 都来自 "mr" 所拥有的内存块。
\end{codebox}

我选择这个例子，是为了让一切都显得简单直白；而且我略去了 \texttt{Chunky\allowbreak{}Ptr<\allowbreak{}T>\allowbreak{}} 类型本身的诸多细节。如果你自己动手抄写这段代码，会发现必须给 \texttt{Chunky\allowbreak{}Ptr} 提供大量重载运算符，比如 \texttt{==}、\texttt{!=}、\texttt{<}、\texttt{++}、\texttt{--} 和 \texttt{-}；此外还得为 \texttt{Chunky\allowbreak{}Ptr<\allowbreak{}void>\allowbreak{}} 提供一个特化，省去重载的 \texttt{operator*}。其中大部分细节，与我们在第 2 章「迭代器与范围（Iterators and Ranges）」中实现自己的迭代器类型时介绍的相同。事实上，每一个「花式指针（fancy pointer）」类型都必须能当作随机访问迭代器来使用——这意味着你必须提供第 2 章「迭代器与范围」末尾列出的五个内嵌类型定义：\texttt{iterator\_\allowbreak{}category}、\texttt{difference\_\allowbreak{}type}、\texttt{value\_\allowbreak{}type}、\texttt{pointer} 和 \texttt{reference}。

最后，如果你想使用某些容器，比如 \texttt{std::\allowbreak{}list} 和 \texttt{std::\allowbreak{}map}，还得实现一个名字出人意料的静态成员函数 \texttt{pointer\_\allowbreak{}to(\allowbreak{}r)\allowbreak{}}：

\begin{codebox}[cpp]
    static ChunkyPtr<T> pointer_to(T &r) noexcept {
      return ChunkyPtr<T>(&r, nullptr);
    }
\end{codebox}

原因在于——你可能还记得第 4 章「容器动物园（The Container Zoo）」的内容——像 \texttt{std::\allowbreak{}list} 这样的少数容器，会把数据存放在节点中，而节点的 \texttt{prev} 和 \texttt{next} 指针既要能指向已分配的节点，也要能指向 \texttt{std::\allowbreak{}list} 对象自身成员数据里包含的节点。有两种显而易见的办法：要么把每个 \texttt{next} 指针都存成「花式指针加裸指针」的某种带标签联合（tagged union，比如第 5 章「词汇类型（Vocabulary Types）」中讲过的 \texttt{std::\allowbreak{}variant}）；要么就想办法把裸指针编码成花式指针。标准库选择了后者。因此，每当你编写一个花式指针类型时，它不仅必须满足分配器对它的全部要求，不仅必须满足随机访问迭代器的要求，还必须能够表示程序地址空间中的任意指针——至少，当你想把这个分配器用于 \texttt{std::\allowbreak{}list} 这类基于节点的容器时是这样。

即便跨过了上述所有难关，你也会发现（截至本书交稿时），libc++ 和 libstdc++ 都无法在任何比 \texttt{std::\allowbreak{}vector} 更复杂的容器中处理花式指针。它们仅仅支持到能配合单一花式指针类型——\texttt{boost::\allowbreak{}interprocess::\allowbreak{}offset\_\allowbreak{}ptr<\allowbreak{}T>\allowbreak{}}，一种不携带元数据的指针——工作的程度。而且标准仍在演进：\texttt{std::\allowbreak{}pmr::\allowbreak{}memory\_\allowbreak{}resource} 是 C++17 才新引入的，截至本文写作时，它仍未由 libc++ 或 libstdc++ 实现。

你可能也注意到了，对于使用花式指针的内存资源，标准库并没有提供任何基类。不过没关系，这东西写起来很容易：

\begin{codebox}[cpp]
    namespace my {
      template<class VoidPtr>
      class fancy_memory_resource {
      public:
        VoidPtr allocate(size_t bytes,
          size_t align = alignof(std::max_align_t)) {
          return do_allocate(bytes, align);
        }
        void deallocate(VoidPtr p, size_t bytes,
          size_t align = alignof(std::max_align_t)) {
          return do_deallocate(p, bytes, align);
        }
        bool is_equal(const fancy_memory_resource& rhs) const noexcept {
          return do_is_equal(rhs);
        }
        virtual ~fancy_memory_resource() = default;
      private:
        virtual VoidPtr do_allocate(size_t bytes, size_t align) = 0;
        virtual void do_deallocate(VoidPtr p, size_t bytes,
          size_t align) = 0;
        virtual bool do_is_equal(const fancy_memory_resource& rhs)
          const noexcept = 0;
      };
      using memory_resource = fancy_memory_resource<void*>;
    } // namespace my
\end{codebox}

标准库没有提供任何使用花式指针的分配器；所有库提供的分配器类型都使用裸指针。

\section{将容器固定到单一内存资源}

标准分配器模型戴的下一顶「帽子」——也就是由 \texttt{std::\allowbreak{}allocator\_\allowbreak{}traits} 掌控的下一个特性——是能把特定的容器对象与特定的堆关联起来。前面我们曾用三个要点来描述这个特性：

\begin{itemize}
\item 将容器对象与某个特定的内存资源（memory resource）关联起来，并确保这种关联是「黏性（sticky）」的——这个容器对象将始终使用给定的堆来进行自身的分配。
\item 将容器值（value）与某个特定的内存资源关联起来，也就是说，该容器可以借助值语义被高效地四处移动，而不会忘记如何释放自己的内容。
\item 在上述两种互斥的行为之间作出选择。
\end{itemize}

我们来看一个例子：用 \texttt{std::\allowbreak{}pmr::\allowbreak{}monotonic\_\allowbreak{}buffer\_\allowbreak{}resource} 充当我们的内存资源，但用一个手写的类类型充当我们的分配器类型。（先让你放心：你确实没有漏掉任何内容——事实上，除了 \texttt{std::\allowbreak{}allocator<\allowbreak{}T>\allowbreak{}}，也就是那个指向由 \texttt{new} 和 \texttt{delete} 管理的全局堆的、微不足道的无状态（stateless）分配器之外，我们到现在仍然没有介绍任何标准库提供的分配器类型。）

\begin{codebox}[cpp]
    template<class T>
    struct WidgetAlloc {
      std::pmr::memory_resource *mr;
      using value_type = T;
      WidgetAlloc(std::pmr::memory_resource *mr) : mr(mr) {}
      template<class U>
      WidgetAlloc(const WidgetAlloc<U>& rhs) : mr(rhs.mr) {}
      T *allocate(size_t n) {
        return (T *)mr->allocate(n * sizeof(T), alignof(T));
      }
      void deallocate(void *p, size_t n) {
        mr->deallocate(p, n * sizeof(T), alignof(T));
      }
    };
    class Widget {
      char buffer[10000];
      std::pmr::monotonic_buffer_resource mr {buffer, sizeof buffer};
      std::vector<int, WidgetAlloc<int>> v {&mr};
      std::list<int, WidgetAlloc<int>> lst {&mr};
    public:
      static void swap_elems(Widget& a, Widget& b) {
        std::swap(a.v, b.v);
      }
    };
\end{codebox}

这里我们的 \texttt{Widget} 是一个典型的面向对象（object-oriented）类类型；我们期望它在整个生命周期里都待在某个固定的内存地址上。接着，为了减少堆碎片化或改善缓存局部性（cache locality），我们在每个 \texttt{Widget} 对象内部放了一大块缓冲区，并让 \texttt{Widget} 把这块缓冲区当作其数据成员 \texttt{v} 和 \texttt{lst} 的后备存储（backing store）。

现在看看 \texttt{Widget::\allowbreak{}swap\_\allowbreak{}elems(\allowbreak{}a,\allowbreak{} b)\allowbreak{}} 函数。它会交换 \texttt{Widget a} 和 \texttt{Widget b} 的 \texttt{v} 数据成员。你可能还记得第 4 章「容器动物园」里讲过，\texttt{std::\allowbreak{}vector} 不过就是指向一块动态分配数组的指针，因此通常库只要交换两个 \texttt{std::\allowbreak{}vector} 实例底层的指针，就能完成交换，而不必搬动任何底层数据——这使得向量的交换成为 $O(1)$ 操作，而非 $O(n)$ 操作。

此外，\texttt{vector} 足够聪明，它知道如果交换指针，就也必须交换分配器——这样，关于「如何释放」的信息，就会跟着那个最终需要被释放的指针一起走。

但在本例中，如果库只是简单地交换指针和分配器，那将是一场灾难！我们会得到一个向量 \texttt{a.\allowbreak{}v}，它底层的数组现在「属于」了 \texttt{b.\allowbreak{}mr}，反之亦然。如果我们销毁了 \texttt{Widget b}，那么下次访问 \texttt{a.\allowbreak{}v} 的元素时，我们就会访问到已经被释放的内存。更糟的是，即便我们再也不访问 \texttt{a.\allowbreak{}v}，当 \texttt{a.\allowbreak{}v} 的析构函数试图调用那个早已消亡的 \texttt{b.\allowbreak{}mr} 的 \texttt{deallocate} 方法时，我们的程序也很可能会崩溃！

所幸，标准库把我们从这个命运中拯救了出来。支持分配器的容器（allocator-aware container）有一项职责，就是在拷贝赋值、移动赋值和交换（swap）时恰当地传播它的分配器。出于历史原因，这件事由 \texttt{allocator\_\allowbreak{}traits} 类模板里一大堆乱七八糟的类型定义来处理；不过，为了正确地使用分配器传播（propagation），你只需要知道以下几件事：

\begin{itemize}
\item 分配器究竟是传播（propagate）自身，还是死死黏附在某个特定的容器上，这是分配器类型本身的属性。如果你想让某个分配器「黏附（stick）」而让另一个传播，就必须把它们做成不同的类型。
\item 当一个分配器具有「黏性」时，它会黏附在某个特定的（古典的、面向对象的）容器对象上。对于非黏性的分配器类型而言原本是 $O(1)$ 的指针交换操作，此时可能退化为 $O(n)$，因为把元素从别的分配器的内存空间「接纳（adopt）」进我们自己的空间，需要在我们自己的空间里为它们分配位置。
\item 「黏性」有一个明确的用武之地（正如我们刚才用 \texttt{Widget} 展示的那样），而「非黏性」的后果可能是灾难性的（同样的，参见 \texttt{Widget}）。因此，\texttt{std::\allowbreak{}allocator\_\allowbreak{}traits} 默认假定一个分配器类型是黏性的，除非它能看出该分配器类型为空，从而几乎可以确定它是无状态（stateless）的。对于空的分配器类型，默认行为实际上就是非黏性。
\item 作为程序员，你基本上总是想要这个默认行为：无状态分配器不妨传播，而有状态（stateful）分配器除了像 \texttt{Widget} 这样需要黏性的场景之外，大概没什么用武之地。
\end{itemize}

\section{使用标准分配器类型}

我们来谈谈标准库提供的分配器类型。\texttt{std::\allowbreak{}allocator<\allowbreak{}T>\allowbreak{}} 是默认的分配器类型；它是每个标准容器模板类型参数的默认值。举例来说，当你在代码里写 \texttt{std::\allowbreak{}vector<\allowbreak{}T>\allowbreak{}} 时，它其实就是 \texttt{std::\allowbreak{}vector<\allowbreak{}T,\allowbreak{} std::\allowbreak{}allocator<\allowbreak{}T>\allowbreak{}>} 这个类型的化名。正如本章前面提到过的，\texttt{std::\allowbreak{}allocator<\allowbreak{}T>\allowbreak{}} 是一个无状态的空类型；它是指向由 \texttt{new} 和 \texttt{delete} 管理的全局堆的一个「句柄（handle）」。

因为 \texttt{std::\allowbreak{}allocator} 是一个无状态类型，\texttt{allocator\_\allowbreak{}traits} 便（正确地）假定它应当是非黏性的。这意味着，像 \texttt{std::\allowbreak{}vector<\allowbreak{}T>\allowbreak{}::\allowbreak{}swap} 和 \texttt{std::\allowbreak{}vector<\allowbreak{}T>\allowbreak{}::\allowbreak{}operator=} 这样的操作，保证是极为高效的指针交换——因为任何一个 \texttt{std::\allowbreak{}vector<\allowbreak{}T,\allowbreak{} std::\allowbreak{}allocator<\allowbreak{}T>\allowbreak{}>} 类型的对象，永远知道如何释放「由任何其他 \texttt{std::\allowbreak{}vector<\allowbreak{}T,\allowbreak{} std::\allowbreak{}allocator<\allowbreak{}T>\allowbreak{}>} 对象最初分配的内存」。

\texttt{std::\allowbreak{}pmr::\allowbreak{}polymorphic\_\allowbreak{}allocator<\allowbreak{}T>\allowbreak{}} 是 C++17 新增的类型。它是有状态（stateful）、非空的类型；它唯一的成员数据是一个指向 \texttt{std::\allowbreak{}pmr::\allowbreak{}memory\_\allowbreak{}resource} 的指针。（事实上，它和我们本章前面示例代码里的 \texttt{Widget\allowbreak{}Alloc} 几乎一模一样！）两个不同的 \texttt{std::\allowbreak{}pmr::\allowbreak{}polymorphic\_\allowbreak{}allocator<\allowbreak{}T>\allowbreak{}} 实例未必可以互换，因为它们的指针可能指向截然不同的 \texttt{memory\_\allowbreak{}resource}；这意味着，一个 \texttt{std::\allowbreak{}vector<\allowbreak{}T,\allowbreak{} std::\allowbreak{}pmr::\allowbreak{}polymorphic\_\allowbreak{}allocator<\allowbreak{}T>\allowbreak{}>} 类型的对象，未必知道如何释放「由其他某个 \texttt{std::\allowbreak{}vector<\allowbreak{}T,\allowbreak{} std::\allowbreak{}pmr::\allowbreak{}polymorphic\_\allowbreak{}allocator<\allowbreak{}T>\allowbreak{}>} 最初分配的内存」。反过来，这又意味着 \texttt{std::\allowbreak{}pmr::\allowbreak{}polymorphic\_\allowbreak{}allocator<\allowbreak{}T>\allowbreak{}} 是一种「黏性」的分配器类型；也就是说，像 \texttt{std::\allowbreak{}vector<\allowbreak{}T,\allowbreak{} std::\allowbreak{}pmr::\allowbreak{}polymorphic\_\allowbreak{}allocator<\allowbreak{}T>\allowbreak{}>::\allowbreak{}operator=} 这样的操作，最终可能会做大量的拷贝。

顺带一提，一遍又一遍地把 \texttt{std::\allowbreak{}vector<\allowbreak{}T,\allowbreak{} std::\allowbreak{}pmr::\allowbreak{}polymorphic\_\allowbreak{}allocator<\allowbreak{}T>\allowbreak{}>} 这个类型名写出来相当繁琐。所幸，标准库的实现者们也得出了同样的结论，于是标准库在 \texttt{std::\allowbreak{}pmr} 命名空间中提供了类型别名：

\begin{codebox}[cpp]
    namespace std::pmr {
      template<class T>
      using vector = std::vector<T,
        polymorphic_allocator<T>>;
      template<class K, class V, class Cmp = std::less<K>>
      using map = std::map<K, V, Cmp,
        polymorphic_allocator<typename std::map<K, V>::value_type>>;
      // 其余容器类似
    } // std::pmr 命名空间结束
\end{codebox}

\subsection{设置默认内存资源}

标准的 \texttt{polymorphic\_\allowbreak{}allocator} 与我们示例中的 \texttt{Widget\allowbreak{}Alloc} 最大的不同在于，\texttt{polymorphic\_\allowbreak{}allocator} 是可默认构造的。默认可构造性（default-constructibility）可以说是一个颇具吸引力的分配器特性；它意味着我们可以写出下面两行中的第二行，而不是第一行：

\begin{codebox}[cpp]
    std::pmr::vector<int> v2({1, 2, 3}, std::pmr::new_delete_resource());
        // 指定一个特定的内存资源
    std::pmr::vector<int> v1 = {1, 2, 3};
        // 使用默认内存资源
\end{codebox}

另一方面，当你看着第二行时，可能会纳闷：「底层的数组到底是在哪里分配的呢？」毕竟，指定一个分配器的主要目的，就是想知道自己的字节究竟从何而来！正因为如此，构造一个标准 \texttt{polymorphic\_\allowbreak{}allocator} 的常规做法，是传入一个指向 \texttt{memory\_\allowbreak{}resource} 的指针——事实上，这个惯用法被认为会如此常见，以至于从 \texttt{std::\allowbreak{}pmr::\allowbreak{}memory\_\allowbreak{}resource*} 到 \texttt{std::\allowbreak{}pmr::\allowbreak{}polymorphic\_\allowbreak{}allocator} 的转换是隐式转换。不过，\texttt{polymorphic\_\allowbreak{}allocator} 确实也带有一个默认的、零参数的构造函数。当你默认构造一个 \texttt{polymorphic\_\allowbreak{}allocator} 时，你会得到指向「默认内存资源（default memory resource）」的一个句柄，而默认情况下它就是 \texttt{new\_\allowbreak{}delete\_\allowbreak{}resource(\allowbreak{})\allowbreak{}}。

不过，这一点是可以改变的！指向默认内存资源的指针，被存放在一个全局的原子（atomic，线程安全）变量里，你可以借助库函数 \texttt{std::\allowbreak{}pmr::\allowbreak{}get\_\allowbreak{}default\_\allowbreak{}resource(\allowbreak{})\allowbreak{}}（返回该指针）和 \texttt{std::\allowbreak{}pmr::\allowbreak{}set\_\allowbreak{}default\_\allowbreak{}resource(\allowbreak{})\allowbreak{}}（为指针赋一个新值，并返回之前的值）来操纵它。

如果你想彻底避开经由 \texttt{new} 和 \texttt{delete} 的堆分配，那么在程序启动时调用 \texttt{std::\allowbreak{}pmr::\allowbreak{}set\_\allowbreak{}default\_\allowbreak{}resource(\allowbreak{}std::\allowbreak{}pmr::\allowbreak{}null\_\allowbreak{}memory\_\allowbreak{}resource(\allowbreak{})\allowbreak{})\allowbreak{}} 或许是个明智之举。当然，你无法阻止程序里的其他部分「行差踏错」，自己去调用 \texttt{set\_\allowbreak{}default\_\allowbreak{}resource}；而且，由于这个全局变量被你程序中的每一个线程共享，如果你试图在程序运行过程中去修改默认内存资源，就可能遇到一些非常诡异的行为。比如说，根本没有办法做到「只为我当前的线程设置默认内存资源」。

此外，调用 \texttt{get\_\allowbreak{}default\_\allowbreak{}resource(\allowbreak{})\allowbreak{}}（比如从 \texttt{polymorphic\_\allowbreak{}allocator} 的默认构造函数中调用）会执行一次原子访问，这通常比「本可避免的原子访问」要稍微慢一点。因此，你的最佳策略是避免使用 \texttt{polymorphic\_\allowbreak{}allocator} 的默认构造函数；始终明确地指出你打算使用哪个内存资源。为了绝对的万无一失，你大可以考虑直接用上面那个 \texttt{Widget\allowbreak{}Alloc} 来代替 \texttt{polymorphic\_\allowbreak{}allocator}；由于没有默认构造函数，\texttt{Widget\allowbreak{}Alloc} 根本不可能被误用。

\section{让容器支持分配器}

讲完了内存资源（堆）与分配器（堆的句柄），我们现在转向这副三脚架的第三条腿：容器类。在每一个支持分配器的容器内部，至少必须发生这四件事：

容器实例必须把一个分配器实例作为成员数据存储起来。

（因此，容器必须把分配器的类型作为模板参数；否则它无法知道要为那个成员变量预留多少空间。）

容器必须提供接受分配器参数的构造函数。容器必须真正使用自己的分配器来分配和释放内存；所有对 \texttt{new} 或 \texttt{delete} 的使用都必须被剔除。

容器的移动构造函数、移动赋值运算符和 \texttt{swap} 函数都必须按照其 \texttt{allocator\_\allowbreak{}traits} 来传播分配器。

下面是一个非常简单的「感知分配器的」容器——一个只容纳单个对象、分配在堆上的容器。它有点像是来自第 6 章「智能指针」的 \texttt{std::\allowbreak{}unique\_\allowbreak{}ptr<\allowbreak{}T>\allowbreak{}} 的感知分配器版本：

\begin{codebox}[cpp]
    template<class T, class A = std::allocator<T>>
    class uniqueish {
      using Traits = std::allocator_traits<A>;
      using FancyPtr = typename Traits::pointer;
      A m_allocator;
      FancyPtr m_ptr = nullptr;
    public:
      using allocator_type = A;
      uniqueish(A a = {}) : m_allocator(a) {
        this->emplace();
      }
      ~uniqueish() {
        clear();
      }
      T& value() { return *m_ptr; }
      const T& value() const { return *m_ptr; }
      template<class... Args>
      void emplace(Args&&... args) {
        clear();
        m_ptr = Traits::allocate(m_allocator, 1);
        try {
          T *raw_ptr = static_cast<T *>(m_ptr);
          Traits::construct(m_allocator, raw_ptr,
              std::forward<Args>(args)...
          );
        } catch (...) {
          Traits::deallocate(m_allocator, m_ptr, 1);
          throw;
        }
      }
      void clear() noexcept {
        if (m_ptr) {
          T *raw_ptr = static_cast<T *>(m_ptr);
          Traits::destroy(m_allocator, raw_ptr);
          Traits::deallocate(m_allocator, m_ptr, 1);
          m_ptr = nullptr;
        }
      }
    };
\end{codebox}

注意，\texttt{unique\_\allowbreak{}ptr} 使用 \texttt{T*} 的地方，我们现在的代码用的是 \texttt{allocator\_\allowbreak{}traits<\allowbreak{}A>\allowbreak{}::\allowbreak{}pointer}；而 \texttt{make\_\allowbreak{}unique} 使用 \texttt{new} 和 \texttt{delete} 的地方，我们现在的代码用的是 \texttt{allocator\_\allowbreak{}traits<\allowbreak{}A>\allowbreak{}::\allowbreak{}allocate}/\texttt{construct} 与 \texttt{allocator\_\allowbreak{}traits<\allowbreak{}A>\allowbreak{}::\allowbreak{}destroy}/\texttt{deallocate} 这一组合拳。我们已经讨论过 \texttt{allocate} 和 \texttt{deallocate} 的用途——它们负责从合适的内存资源获取内存。但那些内存块只是原始的字节；要把一块内存变成一个可用的对象，我们必须在那个地址上构造出一个 \texttt{T} 的实例。为此我们可以使用「placement \texttt{new}」语法；不过在下一节中我们会看到，为什么改用 \texttt{construct} 和 \texttt{destroy} 很重要。

最后，在继续之前，请注意 \texttt{uniqueish} 的析构函数会在尝试释放之前先检查是否存在已分配的内存。这很重要，因为它给了我们一个表示「空对象」的 \texttt{uniqueish} 值——这个值可以在不分配任何内存的情况下构造出来，并且是我们这个类型一个合适的「被移动自」的表示。

现在让我们为我们的类型实现移动操作。我们希望确保，当你从一个 \texttt{uniqueish<\allowbreak{}T>\allowbreak{}} 对象中移出之后，被移出的对象是「空的」。此外，如果左侧对象和右侧对象共享同一个分配器，或者如果分配器类型是「非 sticky 的」，那么我们宁愿完全避免调用 \texttt{T} 的移动构造函数——我们想把已分配指针的所有权从右侧对象转移到左侧对象：

\begin{codebox}[cpp]
    uniqueish(uniqueish&& rhs) : m_allocator(rhs.m_allocator)
    {
      m_ptr = std::exchange(rhs.m_ptr, nullptr);
    }
    uniqueish& operator=(uniqueish&& rhs)
    {
      constexpr bool pocma =
        Traits::propagate_on_container_move_assignment::value;
      if constexpr (pocma) {
        // 我们可以采用新的分配器，因为
        // 我们的分配器类型不是「sticky」的。
        this->clear(); // 使用旧的分配器
        this->m_allocator = rhs.m_allocator;
        this->m_ptr = std::exchange(rhs.m_ptr, nullptr);
      } else if (m_allocator() == rhs.m_allocator()) {
        // 我们的分配器被「stuck」在这个容器上；
        // 但由于它与 rhs 的分配器等价，
        // 我们仍然可以采用 rhs 的内存。
        this->clear();
        this->m_ptr = std::exchange(rhs.m_ptr, nullptr);
      } else {
        // 我们不能传播这个新分配器
        // 因此也就无法采用它的内存。
        if (rhs.m_ptr) {
          this->emplace(std::move(rhs.value()));
          rhs.clear();
        } else {
          this->clear();
        }
      }
      return *this;
    }
\end{codebox}

移动构造函数一如既往地简单。唯一微小的差别是，我们必须记得把我们的 \texttt{m\_\allowbreak{}allocator} 构造成右侧对象分配器的一个副本。我们可以用 \texttt{std::\allowbreak{}move} 来移动分配器，而不是拷贝它，但我觉得在这个例子里这样做并不值得。请记住，分配器只是一个指向实际内存资源的轻量「句柄」，而且很多分配器类型（例如 \texttt{std::\allowbreak{}allocator<\allowbreak{}T>\allowbreak{}}）实际上都是空的。拷贝一个分配器类型应该总是相当廉价的。不过，在这里使用 \texttt{std::\allowbreak{}move} 也无伤大雅。

另一方面，移动赋值运算符就非常复杂了！我们要做的第一件事，是检查我们的分配器类型是否「粘性」的。非粘性由 \texttt{propagate\_\allowbreak{}on\_\allowbreak{}container\_\allowbreak{}move\_\allowbreak{}assignment::\allowbreak{}value} 为 \texttt{true} 来表示，我们把它缩写为「\texttt{pocma}」。（实际上，标准规定 \texttt{propagate\_\allowbreak{}on\_\allowbreak{}container\_\allowbreak{}move\_\allowbreak{}assignment} 应该正好就是 \texttt{std::\allowbreak{}true\_\allowbreak{}type} 类型；GNU 的 libstdc++ 会严格要求你满足这一点。所以在定义你自己的分配器类型时要当心。）如果分配器类型是非粘性的，那么移动赋值最高效的做法就是先销毁我们当前的值（如果有的话）——务必使用我们旧的 \texttt{m\_\allowbreak{}allocator}——然后连同分配器一起采用右侧对象的指针。因为我们是连同指针一起采用了分配器，所以我们可以确信，将来我们知道该如何释放这个指针。

另一方面，如果我们的分配器类型是「粘性的」，那么我们就不能采用右侧对象的分配器。如果我们当前（「粘住」的）分配器实例恰好等于右侧对象的分配器实例，那么我们仍然可以采用右侧对象的指针；我们已经知道如何释放由这个特定分配器实例所分配的指针。

最后，如果我们不能采用右侧对象的分配器实例，而且我们当前的分配器实例又不等于右侧对象的，那么我们就不能采用右侧对象的指针——因为在将来的某个时刻我们将不得不释放那个指针，而释放那个指针的唯一办法就是使用右侧对象的分配器实例，可我们又被禁止采用右侧对象的分配器实例，因为我们自己的实例是「粘住」的。在这种情况下，我们实际上必须用自己的分配器实例分配一个全新的指针，然后通过调用 \texttt{T} 的移动构造函数，把数据从 \texttt{rhs.\allowbreak{}value(\allowbreak{})\allowbreak{}} 拷贝到我们自己的值上。这最后一种情况，是唯一一处我们真正调用了 \texttt{T} 的移动构造函数的情形！

拷贝赋值遵循类似的逻辑来传播右侧分配器实例，只不过它查看的特性是 \texttt{propagate\_\allowbreak{}on\_\allowbreak{}container\_\allowbreak{}copy\_\allowbreak{}assignment}，或者说「\texttt{pocca}」。

\texttt{swap} 尤其有趣，因为它的最终情形（当分配器类型是「粘性的」且分配器实例不相等时）需要额外的分配：

\begin{codebox}[cpp]
    void swap(uniqueish& rhs) noexcept {
      constexpr bool pocs =
        Traits::propagate_on_container_swap::value;
      using std::swap;
      if constexpr (pocs) {
        // 我们可以交换分配器，因为
        // 我们的分配器类型不是「sticky」的。
        swap(this->m_allocator, rhs.m_allocator);
        swap(this->m_ptr, rhs.m_ptr);
      } else if (m_allocator == rhs.m_allocator) {
        // 我们的分配器被「stuck」在这个容器上；
        // 但由于它与 rhs 的分配器等价，
        // 我们仍然可以采用 rhs 的内存，反之亦然。
        swap(this->m_ptr, rhs.m_ptr);
      } else {
        // 双方都无法采用对方的内存，并且
        // 因此总有一方必须进行分配。
        auto temp = std::move(*this);
        *this = std::move(rhs); // 可能抛出异常
        rhs = std::move(temp); // 可能抛出异常
      }
    }
\end{codebox}

在那两行标着「可能抛出异常」的地方，我们调用的都是移动赋值运算符，而它在此时可能会调用 \texttt{emplace}，进而向分配器索要内存。如果底层的内存资源已经耗尽，那么 \texttt{Traits::\allowbreak{}allocate(\allowbreak{}m\_\allowbreak{}allocator,\allowbreak{} 1)\allowbreak{}} 很可能会抛出一个异常——这样一来我们就麻烦了，原因有二。第一，我们已经开始搬移状态并释放旧内存，可能会发现无法「回退」到一个合理的状态。第二，也是更重要的一点，\texttt{swap} 是那些极其原始、极其基础的函数之一，标准库根本不为它的失败做任何准备——例如，\texttt{std::\allowbreak{}swap} 算法（第 3 章「迭代器对算法」）被声明为 \texttt{noexcept}，这意味着它必须成功；它不允许抛出异常。因此，如果分配失败发生在我们 \texttt{noexcept} 的 \texttt{swap} 函数之中，我们就会看到一个 \texttt{bad\_\allowbreak{}alloc} 异常穿过调用栈一路上涌，直到抵达我们 \texttt{noexcept} 的 \texttt{swap} 函数声明处；此时 C++ 运行时会停止栈展开并调用 \texttt{std::\allowbreak{}terminate}，而它（除非程序员通过 \texttt{std::\allowbreak{}set\_\allowbreak{}terminate} 改变了它的行为）将导致我们的程序崩溃。

C++17 标准在规范「交换标准容器类型时应当发生什么」这一点上，比这还要更进一步。首先，标准不再说 \texttt{swap} 期间的分配失败会导致对 \texttt{std::\allowbreak{}terminate} 的调用，而是直接说 \texttt{swap} 期间的分配失败会导致未定义行为（UB）。其次，标准并没有把这种未定义行为局限于分配失败！按照 C++17 标准，仅仅是调用 \texttt{swap} 来交换任意两个分配器不相等的、标准库容器实例，就会导致未定义行为（UB），无论是否真的会遇到分配失败！

事实上，libc++ 利用这个优化机会，为所有标准容器的 \texttt{swap} 函数生成的代码大致如下：

\begin{codebox}[cpp]
    void swap(uniqueish& rhs) noexcept {
      constexpr bool pocs =
        Traits::propagate_on_container_swap::value;
      using std::swap;
      if constexpr (pocs) {
        swap(this->m_allocator, rhs.m_allocator);
      }
      // 甚至不检查我们是否知道如何释放
      // 所采用的指针；直接假定我们可以做到。
      swap(this->m_ptr, rhs.m_ptr);
    }
\end{codebox}

请注意，如果你使用这段代码（就像 libc++ 那样）去 \texttt{swap} 两个分配器不相等的容器，你最终会得到指针与其分配器之间的不匹配，然后你的程序在下次试图用那个不匹配的分配器释放其中一个指针时，很可能会崩溃——甚至更糟。当你在处理 C++17 的「便利」类型（例如 \texttt{std::\allowbreak{}pmr::\allowbreak{}vector}）时，牢记这个陷阱是极其重要的！

\begin{codebox}[cpp]
    char buffer[100];
    auto mr = std::pmr::monotonic_buffer_resource(buffer, 100);
    std::pmr::vector<int> a {1,2,3};
    std::pmr::vector<int> b({4,5,6}, &mr);
    std::swap(a, b);
      // 未定义行为
    a.reserve(a.capacity() + 1);
      // 这一行无疑会崩溃，因为
      // 它试图 delete[] 一个栈上的指针
\end{codebox}

如果你的代码设计允许由不同内存资源支撑的容器彼此交换，那么你必须避开 \texttt{std::\allowbreak{}swap}，转而使用下面这个安全的惯用法：

\begin{codebox}[cpp]
    auto temp = std::move(a); // 没问题
    a = std::move(b); // 没问题
    b = std::move(temp); // 没问题
\end{codebox}

我说「避开 \texttt{std::\allowbreak{}swap}」，意思就是「避开 STL 里任何重排算法」，包括诸如 \texttt{std::\allowbreak{}reverse} 和 \texttt{std::\allowbreak{}sort} 这样的算法。这会是相当浩大的工程，所以我并不建议去尝试！

如果你的代码设计允许由不同内存资源支撑的容器彼此交换，那么说真的，你也许该重新考虑一下你的设计。如果你能把它改成：永远只交换共享同一内存资源的容器，或者你能完全避开有状态的、和/或粘性的分配器，那么你永远都不需要操心这个特定的陷阱。

\section{借助 scoped\_\allowbreak{}allocator\_\allowbreak{}adaptor 向下传播}

在上一节中，我们介绍了 \texttt{std::\allowbreak{}allocator\_\allowbreak{}traits<\allowbreak{}A>\allowbreak{}::\allowbreak{}construct(\allowbreak{}a,\allowbreak{} ptr,\allowbreak{} args...)\allowbreak{}}，并将其描述为 placement-\texttt{new} 语法 \texttt{::\allowbreak{}new (\allowbreak{}(\allowbreak{}void*)\allowbreak{}ptr)\allowbreak{} T(\allowbreak{}args...)\allowbreak{}} 的一个更优替代。现在我们来看看，为什么某个特定分配器的作者会想要赋予它不同的语义。

改变我们自己分配器类型的 \texttt{construct} 语义，一个或许显而易见的办法是：让它对原始类型做平凡默认初始化，而不是零初始化。那段代码大致如下：

\begin{codebox}[cpp]
    template<class T>
    struct my_allocator : std::allocator<T>
    {
      my_allocator() = default;
      template<class U>
      my_allocator(const my_allocator<U>&) {}
      template<class... Args>
      void construct(T *p, Args&&... args) {
        if (sizeof...(Args) == 0) {
          ::new ((void*)p) T;
        } else {
          ::new ((void*)p) T(std::forward<Args>(args)...);
        }
      }
    };
\end{codebox}

现在你可以把 \texttt{std::\allowbreak{}vector<\allowbreak{}int,\allowbreak{} my\_\allowbreak{}allocator<\allowbreak{}int>\allowbreak{}>} 当作一个「类向量」类型来使用，它满足 \texttt{std::\allowbreak{}vector<\allowbreak{}int>\allowbreak{}} 的所有常见不变量，唯一的例外是：当你通过 \texttt{v.\allowbreak{}resize(\allowbreak{}n)\allowbreak{}} 或 \texttt{v.\allowbreak{}emplace\_\allowbreak{}back(\allowbreak{})\allowbreak{}} 隐式创建新元素时，新元素会像栈上的变量一样被创建为未初始化状态，而不是被零初始化。

从某种意义上说，我们在这里设计的是一个「适配器（adaptor）」，它套在 \texttt{std::\allowbreak{}allocator<\allowbreak{}T>\allowbreak{}} 之上，以有趣的方式修改了它的行为。如果我们能以同样的方式修改或「适配」任意分配器，那就更好了；要做到这一点，我们只需把 \texttt{template<\allowbreak{}class T>\allowbreak{}} 改成 \texttt{template<\allowbreak{}class A>\allowbreak{}}，并继承自 \texttt{A}，而旧代码是继承自 \texttt{std::\allowbreak{}allocator<\allowbreak{}T>\allowbreak{}}。当然，我们新适配器的模板参数列表不再以 \texttt{T} 开头，所以我们就得自己实现 \texttt{rebind}；这条路很快会陷入深奥的元编程，所以我就不岔开去展示了。

不过，我们还有另一种有用的方式，可以摆弄我们自己分配器类型的 \texttt{construct} 方法。考虑下面这段创建 \texttt{int} 的向量的向量的代码示例：

\begin{codebox}[cpp]
    std::vector<std::vector<int>> vv;
    vv.emplace_back();
    vv.emplace_back();
    vv[0].push_back(1);
    vv[1].push_back(2);
    vv[1].push_back(3);
\end{codebox}

假设我们想把这个容器「粘」到一个我们自己设计的内存资源上，比如我们最钟爱的 \texttt{Widget\allowbreak{}Alloc}。那我们就得写出类似下面这样重复啰嗦的代码：

\begin{codebox}[cpp]
    char buffer[10000];
    std::pmr::monotonic_buffer_resource mr {buffer, sizeof buffer};
    using InnerAlloc = WidgetAlloc<int>;
    using InnerVector = std::vector<int, InnerAlloc>;
    using OuterAlloc = WidgetAlloc<InnerVector>;
    std::vector<InnerVector, OuterAlloc> vv(&mr);
    vv.emplace_back(&mr);
    vv.emplace_back(&mr);
    vv[0].push_back(1);
    vv[1].push_back(2);
    vv[1].push_back(3);
\end{codebox}

请注意，分配器对象的初始化器 \texttt{\&mr} 在两个层级上都重复出现了。需要重复写出 \texttt{\&mr}，使得我们很难在泛型上下文中使用自己的向量 \texttt{vv}；例如，我们没法轻易把它传给某个函数模板去填充数据，因为每次被调用方想要 \texttt{emplace\_\allowbreak{}back} 一个新的「int 的向量」时，它都得知道只有调用方才知道的地址 \texttt{\&mr}。我们想做的，是把「每当你构造向量之向量的一个元素时，都需要把 \texttt{\&mr} 追加到参数列表末尾」这个概念封装并固化下来。而标准库已经替我们考虑到这一点了！

自 C++11 起，标准库（在名为 \texttt{<scoped\_\allowbreak{}allocator>\allowbreak{}} 的头文件中）提供了一个名为 \texttt{scoped\_\allowbreak{}allocator\_\allowbreak{}adaptor<\allowbreak{}A>\allowbreak{}} 的类模板。就像我们那个默认初始化的「适配器」一样，\texttt{scoped\_\allowbreak{}allocator\_\allowbreak{}adaptor<\allowbreak{}A>\allowbreak{}} 继承自 \texttt{A}，从而继承了 \texttt{A} 的全部行为；然后它重写了 \texttt{construct} 方法，去做一些不同的事情。具体来说，它会尝试判断它当前正在构造的 \texttt{T} 对象是否「使用了一个分配器」，如果是，它就会把自己作为一个额外的参数向下传递，传给 \texttt{T} 的构造函数。

要判断类型 \texttt{T} 是否「使用了一个分配器」，\texttt{scoped\_\allowbreak{}allocator\_\allowbreak{}adaptor<\allowbreak{}A>\allowbreak{}::\allowbreak{}construct} 会求助于类型特性 \texttt{std::\allowbreak{}uses\_\allowbreak{}allocator\_\allowbreak{}v<\allowbreak{}T,\allowbreak{}A>\allowbreak{}}，而它（除非你对其做了特化，而你大概不该这么做）为真，当且仅当 \texttt{A} 能隐式转换为 \texttt{T::\allowbreak{}allocator\_\allowbreak{}type}。如果 \texttt{T} 没有 \texttt{allocator\_\allowbreak{}type}，那么库会假定 \texttt{T} 并不关心分配器，只有 \texttt{pair} 和 \texttt{tuple} 这两种特殊情况除外（它们的构造函数都有专门的重载，目的就是把分配器向下传播给自己的成员），以及 \texttt{promise} 这一特殊情况（它可以用一个分配器来分配自己的共享状态，尽管它事后并没有提供任何方式来引用那个分配器对象；我们

可以说，\texttt{promise} 的分配器支持比我们在第 5 章「词汇类型」中看到的类型擦除（type erasure）示例擦除得还要彻底。

由于历史原因，感知分配器（allocator-aware）类型的构造函数可以遵循两种不同的模式，而 \texttt{scoped\_\allowbreak{}allocator\_\allowbreak{}adaptor} 足够聪明，能够识别这两种模式。

较早且较简单的类型（即除 \texttt{tuple} 和 \texttt{promise} 之外的所有类型）往往具有形如 \texttt{T(\allowbreak{}args...,\allowbreak{} A)\allowbreak{}} 的构造函数，其中分配器 \texttt{A} 放在最后。而对于 \texttt{tuple} 和 \texttt{promise}，标准库引入了一种新模式：

\texttt{T(\allowbreak{}std::\allowbreak{}allocator\_\allowbreak{}arg,\allowbreak{} A,\allowbreak{} args...)\allowbreak{}}，其中分配器 \texttt{A} 放在开头，但在它之前有一个特殊的标签值 \texttt{std::\allowbreak{}allocator\_\allowbreak{}arg}，其唯一用途就是表明参数列表中的下一个参数代表一个分配器——这与标签 \texttt{std::\allowbreak{}nullopt} 的唯一用途是表明某个 \texttt{optional} 没有值（见第 5 章「词汇类型」）类似。正如标准禁止创建 \texttt{std::\allowbreak{}optional<\allowbreak{}std::\allowbreak{}nullopt\_\allowbreak{}t>\allowbreak{}} 类型一样，如果你试图创建 \texttt{std::\allowbreak{}tuple<\allowbreak{}std::\allowbreak{}allocator\_\allowbreak{}arg\_\allowbreak{}t>\allowbreak{}}，也会陷入一堆麻烦之中。

使用 \texttt{scoped\_\allowbreak{}allocator\_\allowbreak{}adaptor}，我们可以把前面那个笨拙的示例改写得稍微不那么笨拙：

\begin{codebox}[cpp]
    char buffer[10000];
    std::pmr::monotonic_buffer_resource mr {buffer, sizeof buffer};
    using InnerAlloc = WidgetAlloc<int>;
    using InnerVector = std::vector<int, InnerAlloc>;
    using OuterAlloc =
    std::scoped_allocator_adaptor<WidgetAlloc<InnerVector>>;
    std::vector<InnerVector, OuterAlloc> vv(&mr);
    vv.emplace_back();
    vv.emplace_back();
    vv[0].push_back(1);
    vv[1].push_back(2);
    vv[1].push_back(3);
\end{codebox}

请注意，分配器类型变得更笨拙了，但重要的是，传给 \texttt{emplace\_\allowbreak{}back} 的 \texttt{\&mr} 参数消失了；我们现在可以在那些期望以自然方式压入元素的上下文中使用 \texttt{vv}，而不必处处都记得加上 \texttt{\&mr}。就我们的情况而言，由于使用的 \texttt{Widget\allowbreak{}Alloc} 不可默认构造，一旦忘记 \texttt{\&mr}，症状就是一大堆编译期错误。但你或许还记得本章前面几节说过，\texttt{std::\allowbreak{}pmr::\allowbreak{}polymorphic\_\allowbreak{}allocator<\allowbreak{}T>\allowbreak{}} 会欣然允许你默认构造它，这可能带来灾难性的后果；所以如果你打算使用 \texttt{polymorphic\_\allowbreak{}allocator}，明智的做法或许也是研究一下 \texttt{scoped\_\allowbreak{}allocator\_\allowbreak{}adaptor}，以便减少那些可能忘记指定分配策略的地方。

\subsection{传播不同的分配器}

在介绍 \texttt{scoped\_\allowbreak{}allocator\_\allowbreak{}adaptor<\allowbreak{}A>\allowbreak{}} 时，我还漏掉了一个复杂之处。模板参数列表并不局限于一个分配器类型参数！你实际上可以创建带有多个分配器类型参数的作用域分配器类型，像下面这样：

\begin{codebox}[cpp]
    using InnerAlloc = WidgetAlloc<int>;
    using InnerVector = std::vector<int, InnerAlloc>;
    using MiddleAlloc = std::scoped_allocator_adaptor<
      WidgetAlloc<InnerVector>,
      WidgetAlloc<int>
    >;
    using MiddleVector = std::vector<InnerVector, MiddleAlloc>;
    using OuterAlloc = std::scoped_allocator_adaptor<
      WidgetAlloc<MiddleVector>,
      WidgetAlloc<InnerVector>,
      WidgetAlloc<int>
    >;
    using OuterVector = std::vector<MiddleVector, OuterAlloc>;
\end{codebox}

设置好这些 typedef 之后，我们继续建立三个不同的内存资源（memory resource），并构造一个 \texttt{scoped\_\allowbreak{}allocator\_\allowbreak{}adaptor} 实例，它能够记住全部三个内存资源（因为它内部包含三个不同的 \texttt{Widget\allowbreak{}Alloc} 实例，每个「层级」各一个）：

\begin{codebox}[cpp]
    char bi[1000];
    std::pmr::monotonic_buffer_resource mri {bi, sizeof bi};
    char bm[1000];
    std::pmr::monotonic_buffer_resource mrm {bm, sizeof bm};
    char bo[1000];
    std::pmr::monotonic_buffer_resource mro {bo, sizeof bo};
    OuterAlloc saa(&mro, &mrm, &mri);
\end{codebox}

最后，我们可以构造一个 \texttt{Outer\allowbreak{}Vector} 实例，传入我们的 \texttt{scoped\_\allowbreak{}allocator\_\allowbreak{}adaptor} 参数；这就够了！我们精心打造的分配器类型深处那个被重写的 \texttt{construct} 方法，会负责把参数 \texttt{\&bm} 或 \texttt{\&bi} 传递给任何需要它们的构造函数：

\begin{codebox}[cpp]
    OuterVector vvv(saa);
    vvv.emplace_back();
      // 这次分配来自缓冲区 "bo"。
    vvv[0].emplace_back();
      // 这次分配来自缓冲区 "bm"。
    vvv[0][0].emplace_back(42);
      // 这次分配来自缓冲区 "bi"。
\end{codebox}

正如你所见，深度嵌套的 \texttt{scoped\_\allowbreak{}allocator\_\allowbreak{}adaptor} 并不适合胆小的人；而且它们只有在你一路上创建大量「辅助」typedef 时才真正可用，就像我们在本例中所做的那样。

关于 \texttt{std::\allowbreak{}scoped\_\allowbreak{}allocator\_\allowbreak{}adaptor<\allowbreak{}A...>} 还有最后一点说明：如果容器的嵌套深度超过了模板参数列表中分配器类型的数量，那么 \texttt{scoped\_\allowbreak{}allocator\_\allowbreak{}adaptor} 会表现得像是参数列表中的最后一个分配器类型永远重复下去。例如：

\begin{codebox}[cpp]
    using InnerAlloc = WidgetAlloc<int>;
    using InnerVector = std::vector<int, InnerAlloc>;
    using MiddleAlloc = std::scoped_allocator_adaptor<
      WidgetAlloc<InnerVector>
    >;
    using MiddleVector = std::vector<InnerVector, MiddleAlloc>;
    using TooShortAlloc = std::scoped_allocator_adaptor<
      WidgetAlloc<MiddleVector>,
      WidgetAlloc<InnerVector>
    >;
    using OuterVector = std::vector<MiddleVector, TooShortAlloc>;
    TooShortAlloc tsa(&mro, WidgetAlloc<InnerVector>(&mri));
    OuterVector tsv(tsa);
    tsv.emplace_back();
      // 这次分配来自缓冲区 "bo"。
    tsv[0].emplace_back();
      // 这次分配来自缓冲区 "bi"。
    tsv[0][0].emplace_back(42);
      // 这次分配再次来自缓冲区 "bi"！
\end{codebox}

我们其实在第一个 \texttt{scoped\_\allowbreak{}allocator\_\allowbreak{}adaptor} 示例（也就是涉及 \texttt{vv} 的那个）中就已经依赖了这种行为，尽管我当时没有提到。既然你现在知道了，不妨回过头去研究那个示例，看看「永远重复」的行为用在了哪里，以及如果你想为 \texttt{int} 的内层数组和 \texttt{Inner\allowbreak{}Vector} 的外层数组使用不同的内存资源，该如何修改那段代码。

\section{小结}

分配器本质上是 C++ 中一个相当晦涩的主题，这主要源于历史原因。几种不同的接口、各自对应着晦涩的使用场景，被一层层堆叠在一起；它们全都涉及大量的元编程；而且即便是像花式指针这样相对古老的 C++11 特性，许多厂商对这些特性的支持仍然欠缺。

C++17 提供了标准库类型 \texttt{std::\allowbreak{}pmr::\allowbreak{}memory\_\allowbreak{}resource} 来厘清内存资源（也叫堆）与分配器（也叫堆的句柄）之间既有的区别。内存资源提供 \texttt{allocate} 和 \texttt{deallocate} 方法；分配器除了提供这些方法，还提供 \texttt{construct} 和 \texttt{destroy}。

如果你要实现自己的分配器类型 \texttt{A}，它必须是一个模板；其第一个模板参数应当是它期望 \texttt{allocate} 的类型 \texttt{T}。你的分配器类型 \texttt{A} 还必须有一个模板化的构造函数，以支持从 \texttt{A<\allowbreak{}U>\allowbreak{}} 到 \texttt{A<\allowbreak{}T>\allowbreak{}} 的「重绑定」。与任何其他种类的指针一样，分配器类型必须支持 \texttt{==} 和 \texttt{!=} 运算符。

堆的 \texttt{deallocate} 方法允许要求传入的指针附带额外的元数据。C++ 通过花式指针来处理这一点。C++17 的 \texttt{std::\allowbreak{}pmr::\allowbreak{}memory\_\allowbreak{}resource} 不支持花式指针，但自己实现起来很容易。

花式指针类型必须满足随机访问迭代器的所有要求，必须可置空（nullable），并且必须能转换为普通裸指针。如果你想把花式指针类型用于 \texttt{std::\allowbreak{}list} 这类基于节点的容器，就必须为它提供一个静态的 \texttt{pointer\_\allowbreak{}to} 成员函数。

C++17 区分「粘性」与「非粘性」分配器类型。无状态的分配器类型（例如 \texttt{std::\allowbreak{}allocator<\allowbreak{}T>\allowbreak{}}）是非粘性的；有状态的分配器类型（例如 \texttt{std::\allowbreak{}pmr::\allowbreak{}polymorphic\_\allowbreak{}allocator<\allowbreak{}T>\allowbreak{}}）默认是粘性的。要创建具有非默认粘性的分配器类型，需要设置全部三个成员 typedef，它们被通俗地称为「POCCA」「POCMA」和「POCS」。像 \texttt{std::\allowbreak{}pmr::\allowbreak{}polymorphic\_\allowbreak{}allocator<\allowbreak{}T>\allowbreak{}} 这样的粘性分配器类型，主要——或许只在——古典面向对象的场景中才有用，即容器对象被固定在某个特定内存地址的情况。以值（value）为导向的编程（伴随着大量的移动和交换）需要无状态分配器类型，否则就要求程序中的每个人都使用同一个堆，以及单一一个粘性但实质上无状态的分配器类型。

\texttt{scoped\_\allowbreak{}allocator\_\allowbreak{}adaptor<\allowbreak{}A...>} 有助于简化使用自定义分配器或内存资源的深度嵌套容器。几乎任何使用非默认分配器类型的深度嵌套容器，都需要大量辅助 typedef 才能勉强保持可读性。

交换两个带有不相等粘性分配器的容器：理论上这会引发未定义行为，实际上它会破坏内存并导致段错误。千万别这么做！
